\SourceLeaf{351}%
% A544-P0351-D01
\[
\begin{array}{rll}
\Theta(X^{00}+Y^{00}i)=&2,&\text{je nachdem }\dfrac{M}{m}\text{ im }1.\\
&3,&\hfill2.\\
&0&\hfill3.\\
&1&\hfill4.\text{ Quadr. liegt, und }k^{00}=1
\end{array}
\]

% A544-P0351-D02
III. Liegt $m$ im dritten Quadranten, so wird für $x=\frac12,y=0$; $\eta=\frac12b$ eine ganze Zahl, wofür $X^0+Y^0i=\frac12A+\frac12Bi$. Man setzt dann
\[
\Theta(X^0+Y^0i)=\Theta(\tfrac12A+(\tfrac12B-1)i)
\]
so oft $\beta$ negativ ist.

% A544-P0351-D03
IV. Liegt $m$ im vierten Quadranten, so ist für $\eta=0$,
$\Theta(X^0+Y^0i=0,1,2,3$ zu setzen, je nachdem $\frac{M}{m}$ im $1.\ 2.\ 3.\ 4.$ Quadranten liegt
$k^0=0$.

% A544-P0351-D04
\Article{14}
Aus den vorhergehenden Untersuchungen folgt nunmehr folgende Bestimmung des Decidenten.

Man sammle alle Werthe von $x$ und $y$, die \emph{innerhalb} der Grenzen $0$ und $\frac12$ liegen und wofür entweder $\eta$ und $X$ oder $\eta$ und $Y$ eine ganze Zahl ist, und bestimme für jedes $x+iy$ nach den Regeln des $12$. Art. den Werth von $k$.

% A544-P0351-D05
Man sammle ferner alle Werthe auf den Grenzen d.\,i. wo entweder $x=0$ oder $\frac12$, während $y$ \emph{zwischen} $0$ und $\frac12$, oder $y=0$ oder $=\frac12$, während $x$ \emph{zwischen} $0$ und $\frac12$, die so beschaffen sind, dass $\eta$ eine ganze Zahl und $[\xi]$ ungerade, und bestimme das zugehörige $l$ auf folgende Weise. Es sei $\Theta M(x+yi)=\pm\theta$, das obere Zeichen für gerade, das untere für ungerade $\eta$
\[
\begin{array}{c|c|c|c|c}
&\multicolumn{4}{c}{\text{so ist für }m\text{ im}}\\
\text{für}&1.\text{ Quadr.}&2.\text{ Quadr.}&3.\text{ Quadr.}&4.\text{ Quadr.}\\
y=0&l=-\theta&l=-\theta&l=+\theta&l=+\theta\\
x=\frac12&l=-\theta&l=+\theta&l=+\theta&l=-\theta\\
y=\frac12&l=+\theta&l=+\theta&l=-\theta&l=-\theta\\
x=0&l=+\theta&l=-\theta&l=-\theta&l=+\theta
\end{array}
\]

\SourceLeaf{352}%
% A544-P0352-D01
Kürzer so
\[
l=\pm\theta,
\]
\[
\begin{array}{ll}
\text{das Zeichen ist dasselbe wie das von }a&\text{wenn }x=0\\
\text{das entgegengesetzte}&\text{wenn }x=\frac12\\
\text{dasselbe wie das von }b&\text{wenn }y=\frac12\\
\text{das entgegengesetzte}&\text{wenn }y=0
\end{array}
\]

% A544-P0352-D02
Zu $\Sigma k+\Sigma l$ kommt dann noch hinzu
\[
\left.\begin{array}{ll}
\text{wenn }m\text{ im zweiten Quadranten liegt:}&2,1,0,3\\
\text{wenn }m\text{ im vierten Quadranten liegt:}&0,
\end{array}\right\}
\begin{array}{l}\text{je nachdem }\frac{M}{m}\text{ im}\\1.\ 2.\ 3.\ 4.\text{ Quadr.}\end{array}
\]
wenn $m$ im ersten Quadranten liegt und $\frac12(a-1)$ ungerade ist
\[
\begin{array}{ll}
\Theta(\frac12A+(\frac12B-1)i)&\text{wenn }\beta\text{ positiv}\\
\Theta(\frac12A+\frac12Bi)&\text{wenn }\beta\text{ negativ}
\end{array}
\]
wenn $m$ im dritten Quadranten liegt und $\frac12(a-1)$ ungerade ist
\[
\begin{array}{ll}
-\Theta(\frac12A+\frac12Bi)&\text{wenn }\beta\text{ positiv}\\
-\Theta(\frac12A+(\frac12B-1)i)&\text{wenn }\beta\text{ negativ.}
\end{array}
\]
\BookRule{.3\textwidth}

% A544-P0352-D03
\par\bigskip\pdfbookmark[0]{Zur Theorie der biquadratischen Reste. IV.}{nachlass-biquadratic-IV}\centerline{\textbf{[IV.]}}\medskip
\renewcommand{\Article}[1]{\par\medskip\Needspace{4\baselineskip}\pdfbookmark[1]{Art. #1}{nachlass-biquadratic-IV-art-#1}\centerline{\textbf{#1.}}\nobreak\smallskip}
\Article{1}
\[
\begin{array}{lll}
\text{Biquadratischer Rest?}&m=a+bi;&aa+bb=d\\[3pt]
\text{Modulus}&M=A+Bi,&AA+BB=D
\end{array}
\]
\[
\frac{mD}{M}=\mu\qquad
\mu=\alpha+\beta i,\quad\alpha=aA+bB,\quad\beta=Ab-Ba
\]
\[
\xi+\eta i=\pi;\quad\pi m=x+yi=p;\quad\pi M=X+Yi=P
\]
Relationen
\[
\begin{array}{lll}
x=a\xi-b\eta&d\xi=ax+by&D\xi=AX+BY\\
y=b\xi+a\eta&d\eta=-bx+ay&D\eta=-BX+AY\\
X=A\xi-B\eta&dX=\alpha x+\beta y&Dx=\alpha X-\beta Y\\
Y=B\xi+A\eta&dY=-\beta x+\alpha y&Dy=\beta X+\alpha Y
\end{array}
\]

\SourceLeaf{353}%
% A544-P0353-D01
\[
\begin{aligned}
\beta\xi&=-Bx+bX=Ay-aY\\
\beta\eta&=-Ax+aX=-By+bY\\
\alpha\xi&=Ax+aY=By+aX\\
\alpha\eta&=-Bx+aY=Ay-bX
\end{aligned}
\]

% A544-P0353-D02
Diejenigen $\pi$, wo $\xi$ und $\eta$ zwischen $0$ und $\frac12$ liegen, sollen durch $\pi^0$ bezeichnet werden, und die entsprechenden $p$ und $P$ durch $p^0$ und $P^0$; diejenigen $\pi$, wo $\eta=0$ und $\xi$ zwischen $0$ und $\frac12$, durch $\pi'$; die, wo $\xi=\frac12$ und $\eta$ zwischen $0$ und $\frac12$, durch $\pi''$; diejenigen $\pi$, wo $\eta=\frac12$ und $\xi$ zwischen $0$ und $\frac12$, durch $\pi'''$; endlich die wo $\xi=0$ und $\eta$ zwischen $0$ und $\frac12$, durch $\pi''''$.

Der Decident von $\frac{m}{M}$ wird so gefunden:

% A544-P0353-D03
Man sammle alle \emph{ganzen} $P^0$, für welche mithin $x^0$ und $y^0$ gebrochen sein werden; die respectiven Intensoren von $p^0$ seien $t^0$ d.\,i. die Zahlen $0,1.\ 2,3$, je nachdem
\[
\begin{array}{lllll}
[x^0]&\text{gerade,}&\text{ungerade,}&\text{ungerade,}&\text{gerade}\\
\relax[y^0]&\text{gerade,}&\text{gerade,}&\text{ungerade,}&\text{ungerade}
\end{array}
\]
So ist der gesuchte Decident $=\Sigma\pm t^0$, wo das obere Zeichen für gerade $P^0$, das untere für die ungeraden zu nehmen ist.

Dies ist die \emph{erste} Methode.

% A544-P0353-D04
\Article{2}
Wir wollen nun die einzelnen $P^0$ nach den Werthen von $Y^0$ zusammenordnen. Indem wir uns auf den Fall einschränken, wo $a,b,A,B$ positiv sind, ist der kleinste Werth von $Y^0\dotstwo+1$, der grösste $\frac12(A+B-1)$. Für jeden bestimmten Werth von $Y^0$ müssen die Werthe von $X^0$ zwischen bestimmten Grenzen liegen, nemlich

I. wenn $A-B$ positiv ist
\[
\begin{array}{c|c|c}
\text{wenn}&\text{zwischen}&\text{und}\\
Y<\frac12B&-\frac{BY^0}{A}&\frac{AY^0}{B}\\[5pt]
Y=\frac12B&-\frac{BB}{2A}&\frac12A\\[5pt]
Y>\frac12B\text{ und }<\frac12A&-\frac{BY^0}{A}&-\frac{BY^0}{A}+\frac{D}{2A}\\[5pt]
Y>\frac12A&\frac{AY^0}{B}-\frac{D}{2B^0}&-\frac{BY^0}{A}+\frac{D}{2A}
\end{array}
\]

\SourceLeaf{354}%
% A544-P0354-D01
II. Wenn $A-B$ negativ ist,
\[
\begin{array}{c|c|c}
\text{wenn}&\text{zwischen}&\text{und}\\
Y<\frac12A&-\frac{BY^0}{A}&\frac{AY^0}{B}\\[5pt]
Y>\frac12A\text{ und }<\frac12B&\frac{AY^0}{B}-\frac{D}{2B}&\frac{AY^0}{B}\\[5pt]
Y=\frac12B&\frac{AB-D}{2B}&\frac12A\\[5pt]
Y>\frac12B&\frac{AY^0}{B}-\frac{D}{2B}&-\frac{BY^0}{A}+\frac{D}{2A}
\end{array}
\]

% A544-P0354-D02
In den kleinern Grenzen ist entweder $\xi=0$ oder $\eta=\frac12$, in den grössern Grenzen hingegen entweder $\eta=0$ oder $\xi=\frac12$. Es lässt sich leicht beweisen, dass nie die Grenzen von $x$ ganze Zahlen sind.

% A544-P0354-D03
\Article{3}
Wir wollen nun die Partialsummen für jedes bestimmte $Y^0$ auf eine andere Weise darstellen. Auf den Grenzen wird $p$ bestimmte Werthe haben, die durch $p^*,p^{**}$ bezeichnet werden mögen, und während $X$ stetig von der einen Grenze zur andern sich ändert, wird $p$ stetig von $p^*$ zu $p^{**}$ übergehen. Allein die in $[p]$ enthaltene ganze Zahl wird hiebei sprungsweise geändert, indem immer entweder der reelle oder der imaginäre Theil sich um Eine Einheit ändert. Es geschehen die Aenderungen bei den Werthen von $X$
\[
X',\ X'',\ X'''\ldots X^\mu
\]
die bereits nach ihrer Grösse geordnet sind und denen die Werthe von $p$
\[
p',\ p''.\ p'''\ldots
\]
entsprechen.

% A544-P0354-D04
Das letzte $X^\mu$ kann auch mit $X^{**}$ identisch sein, wenn $\beta$ positiv, oder $X'$ mit $X^*$ identisch etc.

Die $x$ sind hier zunehmend, also wenn $x^{**}$ eine ganze Zahl, wird sie für $x^\mu$ gezählt.

Die $y$ sind zunehmend bei positiven $\beta$, da wird $y^{**}$ ganz mitgezählt
\begin{center}abnehmend bei negativen $\beta$, da wird $y^*$ mitgezählt.\end{center}

\SourceLeaf{355}%
% A544-P0355-D01
Die Intensoren von $p^*,p^{**}$ seien $\lambda^*$ und $\lambda^{**}$
\[
\begin{array}{rll}
\text{der Intensor von }p'&\text{an bis }p''\dotstwo&\lambda^*+\delta'\\
p''&\text{bis }p'''\dotstwo&\lambda^*+\delta'+\delta''\\
p^\mu&\text{bis }p^{**}\dotstwo&\lambda^*+\delta'+\delta''\dotstwo+\delta^\mu=\lambda^{**}
\end{array}
\]
so wird die Partialsumme, in sofern $Y^0$ gerade,
\[
\begin{aligned}
={}&\lambda^*(g-h)+(\lambda^*+\delta')(g'-h')\\
&+(\lambda^*+\delta'+\delta'')(g''-h'')+\text{etc.}+\lambda^{**}(g^\mu-h^\mu)
\end{aligned}
\]
wo $g$ die Anzahl der geraden $X^0$ von $X^*$ bis $X'$, $h$ die der ungeraden bedeutet.

% A544-P0355-D02
\Article{4}
Diese Formel lässt sich auch so darstellen:
\[
\begin{aligned}
&\lambda^*\{[\tfrac12X^*-\tfrac12]-[\tfrac12X^*]\}\\
&+\delta'\{[\tfrac12X'-\tfrac12]-[\tfrac12X']\}\\
&+\delta''\{[\tfrac12X''-\tfrac12]-[\tfrac12X'']\}\\
&+\text{etc.}\\
&+\delta^\mu\{[\tfrac12X^\mu-\tfrac12]-[\tfrac12X^\mu]\}\\
&-\lambda^{**}\{[\tfrac12X^{**}-\tfrac12]-[\tfrac12X^{**}]\}
\end{aligned}
\]
oder durch
\[
\lambda^*\varepsilon^*+\delta'\varepsilon'+\delta''\varepsilon''+\text{etc.}+\delta^\mu\varepsilon^\mu-\lambda^{**}\varepsilon^{**}
\]
\[
\begin{array}{rll}
\text{wo allgemein }\varepsilon=0&\text{wenn }[X]\text{ ungerade}\\
\text{und }=-1&\text{wenn }[X]\text{ gerade ist und }Y\text{ gerade}\\
+1&\text{wenn }[X]\text{ gerade und }Y\text{ ungerade.}
\end{array}
\]

% A544-P0355-D03
Für $\delta$ hingegen hat man die Werthe
\[
\begin{array}{ll|rr}
&&\multicolumn{2}{c}{\beta}\\
&&\text{positiv}&\text{negativ}\\
\text{wenn }x\text{ eine ganze gerade Zahl,}&[y]\text{ gerade}&-1&-1\\
&[y]\text{ ungerade}&+1&+1\\
x\text{ eine ganze ungerade Zahl,}&[y]\text{ gerade}&+1&+1\\
&[y]\text{ ungerade}&-1&-1\\
y\text{ eine ganze gerade Zahl,}&[x]\text{ gerade}&-3&+3\\
&[x]\text{ ungerade}&-1&+1\\
y\text{ eine ganze ungerade Zahl,}&[x]\text{ gerade}&+3&-3\\
&[x]\text{ ungerade}&+1&-1
\end{array}
\]

